% Generated draft: requires mathematical and source review.
\documentclass[11pt]{article}
\usepackage[margin=25mm]{geometry}
\usepackage{fontspec}
\setmainfont{Latin Modern Roman}
\usepackage{amsmath,amssymb,mathtools,unicode-math}
\setmathfont{Latin Modern Math}
\usepackage{enumitem,needspace,xurl,hyperref}
\setlength{\parindent}{0pt}
\setlength{\parskip}{5pt}
\setlength{\emergencystretch}{3em}
\newcommand{\N}{\mathbb{N}}
\newcommand{\Z}{\mathbb{Z}}
\newcommand{\Q}{\mathbb{Q}}
\newcommand{\R}{\mathbb{R}}
\newcommand{\C}{\mathbb{C}}
\newcommand{\F}{\mathbb{F}}
\newcommand{\A}{\mathbb{A}}
\newcommand{\PP}{\mathbb{P}}
\newcommand{\E}{\mathbb{E}}
\newcommand{\eps}{\varepsilon}
\newcommand{\dd}{\,\mathrm{d}}
\begin{document}
\section*{UnsolvedMath ID: 700690}
\textbf{Source problem code:} catalogue entry
\subsection{Definitions and mathematical statement}
Let \(G\) be a group, let \(p\) be a prime, and let \(a\in G\) have order \(p\). Write
\[
a^g=g^{-1}ag,\qquad H_g=\langle a,a^g\rangle,\qquad \mathcal H=\{H_g:g\in G\}.
\]
For a subgroup \(F\leq G\), its normalizer is
\[
N_G(F)=\{g\in G:g^{-1}Fg=F\}.
\]
A group is soluble if its derived series reaches the trivial subgroup after finitely many steps.

The source asks: if the normalizer of every finite subgroup of \(G\) containing \(a\) has finite periodic part, every \(H_g\) is finite, and almost all of these subgroups are soluble, must \(G\) possess a periodic part when \(p>2\)?

The supplied record does not define ``periodic part'' or specify the indexing convention for ``almost all.'' To make the possible interpretation explicit, put
\[
T(K)=\{x\in K:\text{the order of }x\text{ is finite}\}
\]
for any group \(K\). Under the interpretation that possessing a periodic part means that \(T(K)\) is a subgroup, having finite periodic part means that this subgroup is finite, and ``almost all'' means all but finitely many distinct members of \(\mathcal H\), the requested assertion becomes the following conditional statement:

For every group \(G\), every odd prime \(p\), and every element \(a\in G\) of order \(p\), the assumptions
\[
\begin{gathered}
\text{for every finite }F\leq G\text{ with }a\in F,
\\
T(N_G(F))\leq N_G(F)\quad\text{and}\quad |T(N_G(F))|<\infty,
\\[2pt]
|\langle a,a^g\rangle|<\infty\quad\text{for every }g\in G,
\\[2pt]
\bigl|\{H\in\mathcal H:H\text{ is not soluble}\}\bigr|<\infty
\end{gathered}
\]
imply \(T(G)\leq G\). This explicit formulation is conditional on the stated interpretations and requires comparison with the source terminology. If instead ``almost all'' counts elements \(g\in G\), its corresponding assumption would be
\[
\bigl|\{g\in G:H_g\text{ is not soluble}\}\bigr|<\infty;
\]
the supplied record does not resolve this distinction.

The source additionally reports that Shunkov's cited 1986 works establish an affirmative answer when \(a\) is a \emph{point}, and exhibit a group satisfying the stated hypotheses with a point \(a\) of order \(2\) but without a periodic part. The definition of a point is referred to the cited 1983 paper and is not supplied here; no definition or additional hypothesis concerning points is reconstructed in this draft. These are supplied historical reports, not independently verified results.
\subsection{Short English statement}
Shunkov asks whether a group containing an element of odd prime order must have a periodic part when the normalizers of all finite subgroups containing that element have finite periodic part, all subgroups generated by the element and a conjugate are finite, and almost all of those subgroups are soluble. The meanings of ``periodic part'' and ``almost all'' require source confirmation. The supplied record labels the problem open and reports an additional positive result for points and a counterexample at prime \(2\).
\subsection{Sources}
\begin{itemize}
\item Supplied local UnsolvedMath registry record, permanent ID \(700690\), exact source problem code \texttt{catalogue entry}, title \emph{Kourovka Notebook Problem 10.75}. The supplied status is \texttt{open}; it is a registry label, not an independently verified assessment. The record gives proposer V. P. Shunkov and proposed year 1986; upstream revision \texttt{d371ee917ac9647fbb06cd6fc529deec6f3dda82}.
\item \emph{The Kourovka Notebook (Unsolved Problems in Group Theory)}, original problem 10.75, 10th issue (1986); the supplied background identifies the 21st issue edition (2026), arXiv:\texttt{1401.0300v48}. The collection description names E. I. Khukhro and V. D. Mazurov as editors. This attribution and edition information are supplied source data.
\item V. P. Shunkov, \emph{Groups with involutions}, Preprints nos. 4, 5, 12 of the Comput. centre of SO AN SSSR, Krasnoyarsk, 1986 (Russian). Cited in the supplied statement for the reported affirmative result under the point assumption and the reported order-\(2\) counterexample.
\item V. I. Senashov and V. P. Shunkov, \emph{Algebra and Logic}, volume 22, no. 1 (1983), pp. 66--81. The supplied statement refers here for the definition of a point; no article title or definition is supplied.
\end{itemize}
\end{document}
